\section{Final Years, Death, and Ecclesial Reception}

Ambrose's last seven years unfolded against the rapid collapse of Western political stability. On 15~May 392, the twenty-year-old Valentinian~II was found dead at Vienne in Gaul---almost certainly murdered at the instigation of his Frankish general Arbogast---while waiting for Ambrose to cross the Alps and baptize him. When Valentinian's body reached Milan, Ambrose pronounced the funeral consolation (\work{De Obitu Valentiniani Consolatio}, summarized in de Romestin's Prolegomena), comforting the boy's grieving sisters with the classic theological teaching on baptism of desire for a catechumen cut off before the font.

\subsection{Eugenius, Frigidus, and the Plea for Clemency (392--395)}

Unable as a Frank to claim the purple himself, Arbogast elevated the Christian rhetorician Eugenius as Western emperor. Seeking support among Rome's traditionalist aristocracy, Eugenius---after twice refusing---granted revenues to pagan senators that financed the revival of traditional temple rites (\work{Letter} LVII.6). Unwilling to greet a ruler who compromised with idolatry, Ambrose withdrew from Milan in 393 to Bononia (Bologna), Faenza, and Florence, where he consecrated a basilica with the relics of the martyrs Vitalis and Agricola (\work{Prolegomena}~\S III, under 393--394). From exile he wrote \work{Letter} LVII to Eugenius, calmly rehearsing his consistency under Valentinian~II and Theodosius and refusing flattery: ``in the cause of God I felt a righteous fear, and \ldots\ I did not value flattery above my own soul'' (\work{Letter} LVII.12).

Returning to Milan around 1~August 394 as Eugenius marched east to face Theodosius, Ambrose received news of Theodosius's decisive victory at the River Frigidus (6~September 394). In \work{Letter} LXI, Ambrose told the victorious emperor how he had honored his letter of thanksgiving:

\begin{quote}
I took the letter of your Piety with me to the altar. I laid it upon the altar. I held it in my hand whilst I offered the Sacrifice; so that your faith might speak by my voice, and the Emperor's letter discharge the function of the priestly oblation. (\work{Letter} LXI.5)
\end{quote}

Yet in the very same breath, and again through his deacon Felix in \work{Letter} LXII, Ambrose pleaded urgently for mercy toward Eugenius's defeated partisans who had fled into church sanctuary: ``by your clemency the Church of God, as it delights in the peace and tranquillity of the innocent, so, too, may rejoice in the pardon of the guilty'' (\work{Letter} LXI.7); ``I have been unable to endure their tears without anticipating by my entreaty the coming of your Clemency'' (\work{Letter} LXII.3). Four months later, in January 395, Theodosius died in Milan, commending his young sons Honorius and Arcadius to the Church; Ambrose delivered his funeral oration (\work{De Obitu Theodosii}) forty days later.

\subsection{Last Labors and Death (396--397)}

Ambrose's metropolitan care never slackened in old age. In 396, when factional strife kept the neighboring see of Vercellae long vacant after the death of Bishop Limenius, Ambrose wrote his long, magisterial \work{Epistle} LXIII, invoking the memory of the exiled confessor-bishops Eusebius of Vercellae and Dionysius of Milan, rebuking two former Milanese monks---Sarmatio and Barbatianus---who had turned against fasting and continence, and traveling to Vercellae to restore concord (\work{Epistle} LXIII.1--9, 68--70; \work{Prolegomena}~\S III, under 396). Early in 397, he consecrated a bishop for Ticinum (Pavia) and soon afterward fell gravely ill.

The last weeks of his life are preserved chiefly through Paulinus of Milan, who attended him as secretary and wrote his \work{Vita} at Augustine's request (summarized in de Romestin's Prolegomena~\S III and \S V). Even from his sickbed, Ambrose continued dictating his \work{Enarratio} on Psalm~43 (44), leaving it unfinished at his death. When the Western regent Stilicho urged Milan's nobles to beg the bishop to pray for longer life---fearing that Ambrose's passing threatened ruin to Italy---Ambrose answered: ``I have not so lived amongst you as to be ashamed of living, and I do not fear to die, for we have a good Lord.'' Overhearing attendants whisper names of possible successors and mention his aged mentor Simplicianus, he murmured, ``An old man, but good.'' Lying for his final hours with his arms outstretched in the form of a cross, he received the Eucharist as Viaticum from Honoratus, bishop of Vercellae, and breathed his last on the night between Good Friday and Holy Saturday, 4~April 397. On Easter Sunday, 5~April 397, he was laid to rest in the Basilica Ambrosiana beside Saints Gervasius and Protasius, fulfilling the wish he had written to Marcellina eleven years before (\work{Letter} XXII.13).

\subsection{Liturgical Memory and the Doctorate of the Church}

Veneration of Ambrose began immediately at his tomb in Milan and spread rapidly across East and West. Because his death fell on 4~April---often overlapping Holy Week or Easter---both Milan and Rome anchored his principal liturgical feast on 7~December, the anniversary of his episcopal ordination. The eleventh-century Milanese \work{Manuale Ambrosianum} (edited by Marco Magistretti in 1904) preserves the rich medieval Ambrosian office for \latin{Die VII Men.\ Decembris: Ordinatio Sancti Ambrosii}, opening Vespers with the hymn \latin{Miraculum laudabile} (``Sing, all peoples, the praiseworthy miracle''), which celebrates ``Ambrose the great seer, predestined by God'' to defeat Arianism and echoes the popular cry \latin{Ambrosium episcopum} (p.~18), while the morning antiphon (\latin{Ant.\ xxiii}, pp.~21--22) sings of the child's voice proclaiming him bishop: \latin{Vox infantis sonuit: Dignus Ambrosius, hic erit episcopus} (\work{Manuale Ambrosianum}, Pars Altera, pp.~17--22).

Centuries later, the 1640 \work{Missale Ambrosianum} issued by Archbishop Cesare Cardinal Monti of Milan offered a remarkably balanced historical tribute in its preface to the Milanese clergy. Rather than claiming that Ambrose single-handedly invented Milan's liturgy, Cardinal Monti traced its cumulative growth across a line of pastors---from local tradition's memory of St.~Barnabas through Bishop Myrocles to Ambrose himself, ``our Father and Patron,'' who ``not only retained the rites already instituted, but increased them with such profit and fruitfulness that not long after his time our Church received the distinguished name of Ambrosian,'' followed by Simplicianus, who ``added much,'' and Eusebius, who restored the service books after wartime fires (\work{Missale Ambrosianum}, 1640, preface).

In the universal Roman Rite, the 1902 Propaganda Fide \work{Martyrologium Romanum} commemorates Ambrose twice: on 4~April (\latin{Pridie Nonas Aprilis}), recording his deposition at Milan and crediting his teaching and miracles with restoring nearly all Italy from Arianism to the Catholic faith; and on 7~December (\latin{Septimo Idus Decembris}), commemorating the ordination of ``Saint Ambrose, Bishop and Doctor of the Church, by whose sanctity and doctrine the universal Church is adorned.'' The same Roman Martyrology honors his brother Satyrus on 17~September and Saints Gervasius and Protasius on 19~June. Both the 1962 \work{Missale Romanum} (as a Third Class feast on 7~December, p.~435, no.~1863; Latin Collect \latin{Deus, qui populo tuo}, printed alongside English in the 1945 \work{New Roman Missal}, pp.~845--846) and the postconciliar General Roman Calendar (as an Obligatory Memorial) preserve his feast on 7~December, with the traditional Roman Collect praying that ``whom we had as a teacher of life on earth, we may merit to have as an intercessor in heaven'' (\latin{quem Doctorem vitae habuimus in terris, intercessorem habere mereamur in caelis}).

In 1298, Pope Boniface~VIII formally ordered the solemn observance of the feasts of the four original Latin Doctors of the Church---Ambrose, Jerome, Augustine, and Gregory the Great---confirming in canon law a standing that Catholic theology had recognized since the fifth century. Alongside his pupil Augustine, Ambrose remains a foundational doctor of Trinitarian orthodoxy, sacramental realism, Marian faith, clerical ethics, and liturgical praise.
